\section{Giới thiệu}\label{sec:intro}

Trong mô hình ReAct~\cite{yao2023react}, tác tử LLM xen kẽ suy luận với lời gọi công cụ: tìm tài liệu, truy vấn
cơ sở dữ liệu, gửi e-mail hay gọi webhook. Khi đó nội dung công cụ trả về đi thẳng vào ngữ cảnh của mô hình. Trong
tấn công chèn lệnh gián tiếp (IPI)~\cite{greshake2023indirect}, kẻ tấn công không nói chuyện với mô hình mà cài chỉ
thị vào nội dung bên thứ ba mà tác tử sẽ đọc. Nếu mô hình làm theo, nó có thể gửi dữ liệu tới đích lạ, làm lộ giá
trị mật hoặc lặng lẽ sửa dữ liệu người dùng yêu cầu gửi.

Các benchmark tác tử đã công bố như BIPIA~\cite{yi2025bipia}, InjecAgent~\cite{zhan2024injecagent} và
AgentDojo~\cite{debenedetti2024agentdojo} gần như chỉ bằng tiếng Anh. Tiếng Việt có những đặc điểm có thể ảnh hưởng
tới cả tấn công lẫn phòng thủ: văn bản thường bị bỏ dấu, thuật ngữ tiếng Anh hay trộn vào câu, và bộ phân loại chèn
lệnh chủ yếu được huấn luyện trên tiếng Anh. Mặt khác, phòng thủ ở mức prompt như
spotlighting~\cite{hines2024spotlighting} phụ thuộc vào việc mô hình có tuân theo hay không, còn bộ phân loại tốn
thêm chi phí suy luận và chuyển giao kém giữa các ngôn ngữ. Điều đó dẫn tới hướng cưỡng chế \emph{không trông vào
sự tuân thủ của mô hình}: luật do host áp dụng giữa mô hình và công cụ.

Bài báo có bốn đóng góp:
\begin{enumerate}
\item \textbf{Mô hình đe dọa và phân loại theo cơ chế:} hai nhãn độc lập là độ tin cậy và độ nhạy cảm, bốn vị từ
  gây hại, sáu cơ chế tấn công và hai nhánh lành tính làm lộ chặn nhầm (Mục~\ref{sec:threat}).
\item \textbf{\ours}, phòng thủ nhiều lớp phía broker gồm bốn luật tất định cùng lời nhắc và phản hồi từ chối;
  không cần tinh chỉnh và không cần mô hình phụ (Mục~\ref{sec:defence}).
\item \textbf{VN-AgentInject v3 và AgentDojo-VI} với giao thức niêm phong: giả thuyết đăng ký trước, cấu hình ghim
  bằng hash, mỗi tập kiểm thử chạy một lần (Mục~\ref{sec:protocol}). \todo{owner quyết: công bố dữ liệu / cung
  cấp khi có yêu cầu}
\item \textbf{Kết quả:} trên hai tập niêm phong và ba mô hình nền mở trọng số, \ours{} giảm tỉ lệ tấn công thành công
  từ 32,4\% xuống 3,6\% với cái giá 3,4 điểm tiện ích; phát hiện học trên tiếng Anh không chuyển giao sang tiếng
  Việt; lỗi còn lại tập trung ở rò rỉ giá trị mật và sửa tham số (Mục~\ref{sec:results}--\ref{sec:discussion}).
\end{enumerate}
